\clearpage
\section{Historial de Cambios - La evolución de nuestro documento}\label{sec:spmp-1}
\begingroup\small\setlength{\tabcolsep}{5pt}\renewcommand{\arraystretch}{1.2}
\begin{longtable}{|p{0.12\linewidth}|p{0.17\linewidth}|p{\dimexpr0.71\linewidth-6\tabcolsep-4\arrayrulewidth\relax}|}\hline
Fecha & Autor & Descripción\\[2mm]\hline
26/10/2026 & Daniel Ramírez & \vspace{-\baselineskip}\begin{itemize}\setlength{\itemsep}{0pt}\setlength{\parskip}{0pt}\setlength{\topsep}{0pt} \item Creación de la estructura y redacción inicial del PMP con portada, índice y secciones de la plantilla SPMP, a partir del contexto del proyecto, la adaptación de estándares, la documentación del servidor y el canvas del equipo.\end{itemize}\\\hline
27/10/2026 & Daniel Felipe Ramirez Vargas & \vspace{-\baselineskip}\begin{itemize}\setlength{\itemsep}{0pt}\setlength{\parskip}{0pt}\setlength{\topsep}{0pt} \item Actualización de portada y formato \item revisión integral de contenido, estructura de títulos, referencias y terminología FDS/WFDS.\end{itemize}\\\hline
28/10/2026 & Daniel Felipe Ramirez Vargas & \vspace{-\baselineskip}\begin{itemize}\setlength{\itemsep}{0pt}\setlength{\parskip}{0pt}\setlength{\topsep}{0pt} \item Revisión de contenido, actividades, responsabilidades, formación, infraestructura y referencias del PMP.\end{itemize}\\\hline
29/10/2026 & Daniel Felipe Ramirez Vargas & \vspace{-\baselineskip}\begin{itemize}\setlength{\itemsep}{0pt}\setlength{\parskip}{0pt}\setlength{\topsep}{0pt} \item Actualización de los títulos de las secciones.\end{itemize}\\\hline
06/11/2026 & Daniel Felipe Ramirez Vargas & \vspace{-\baselineskip}\begin{itemize}\setlength{\itemsep}{0pt}\setlength{\parskip}{0pt}\setlength{\topsep}{0pt} \item Resolución de pendientes, actualización de objetivos y fuente de igniciones de la sección 6.2.\end{itemize}\\\hline
07/11/2026 & Daniel Felipe Ramirez Vargas & \vspace{-\baselineskip}\begin{itemize}\setlength{\itemsep}{0pt}\setlength{\parskip}{0pt}\setlength{\topsep}{0pt} \item Sección 6.3: resolución de pendientes sobre datos y acceso a HPC \item sección 6.4: pendientes de entregables \item secciones 2, 6.4, 6.5, 7, 8.1, 8.3, 10.3 y 11: denominación Semi-transición \item sección 6.4: referencia a actividades \item \textcolor{black}{sección 6.5: hitos, ítems de presupuesto consolidados y actualización de pendientes} \textcolor{black}{(fechas de H1 y H2)} \item \textcolor{black}{sección 6.6: responsables, aprobación y control de versiones} \item sección 9.3.1: identificadores, fases, relación con entregables y formato de tabla \item formato general de tablas \item sección 7: ampliación del glosario \item secciones 9.2.1 y 9.2.2: capacitación e infraestructura.\end{itemize}\\\hline

08/11/2026 & Miguel Ángel Sarmiento Cuervo &
\vspace{-\baselineskip}
\begin{itemize}
    \setlength{\itemsep}{0pt}
    \setlength{\parskip}{0pt}
    \setlength{\topsep}{0pt}
    \item Sección 11: ampliación del procedimiento de entrega, consolidación, verificación, empaquetado y presentación de los entregables E1--E7.
    \item Sección 12.1: revisión del ambiente de trabajo y los acuerdos del equipo.
    \item Sección 12.2: desarrollo del plan de administración de riesgos, matriz de probabilidad e impacto, acciones preventivas y diagrama BPMN.
    \item Sección 12.3: ampliación de los procedimientos de configuración, control de versiones, documentación, almacenamiento y recuperación.
    \item Sección 12.4: revisión y ampliación del proceso de medición y seguimiento.
    \item Sección 12.5: desarrollo de siete procesos de control de calidad, tabla de responsabilidades y diagramas BPMN.
    \item Revisión de coherencia, compilación y corrección de problemas de diagramación de las secciones 11 y 12.
\end{itemize}\\\hline

08/11/2026 & Andrés Felipe Pinzón Márquez &
\vspace{-\baselineskip}
\begin{itemize}
    \setlength{\itemsep}{0pt}
    \setlength{\parskip}{0pt}
    \setlength{\topsep}{0pt}
    \item Sección 10: desarrollo, redacción y consolidación integral de la sección de Monitoreo y Control del Proyecto.
    \item Sección 10.1: formulación del procedimiento de administración de requerimientos y estructuración del diagrama BPMN de control de cambios.
    \item Sección 10.2: definición de la base de medición, métricas, indicadores de seguimiento y procedimientos ante desviaciones.
    \item Sección 10.3: establecimiento del cierre del proyecto, fases de evaluación, correcciones y entrega con fechas por confirmar y acuerdos colectivos.
\end{itemize}\\\hline

\end{longtable}\endgroup

